%%%%%%%%%%  scan idx 238  |  folio 225  %%%%%%%%%%
\textbf{Lemme 1.1.4.}
\uline{Soit un diagramme commutatif de schémas}
\[
\begin{tikzcd}
X' \arrow[dr,"f'"'] & X \arrow[l,hook',"i'"'] \arrow[r,hook,"i''"] & X'' \arrow[dl,"f''"]\\
& Y{,} &
\end{tikzcd}
\]
\uline{où} \(i'\) (\uline{resp.} \(i''\)) \uline{est une immersion fermée}, \uline{et} \(f'\) (\uline{resp.} \(f''\)) \uline{un morphisme lisse}, \uline{et soit} \(E\in\ob D(X)\). \uline{On a alors} :
\[
  i'_*(E)\ n\text{-\uline{pseudo-cohérent}}
  \quad\Longleftrightarrow\quad
  i''_*(E)\ n\text{-\uline{pseudo-cohérent}}.
\]

\uline{Preuve}. Suivant une astuce due à Lichtenbaum, on envoie \(X\) diagonalement dans \(X'\times_YX''\), et l'on est aussitôt ramené à prouver l'assertion dans le cas où l'un des morphismes \(f'\), \(f''\) est l'identité. Changeant les notations, on a donc un triangle commutatif
\begin{equation}\tag{$*$}
\begin{tikzcd}
& Y \arrow[d,"f"]\\
X \arrow[r,hook,"j"'] \arrow[ur,hook,"i"] & Z{,}
\end{tikzcd}
\end{equation}
où \(i\) et \(j\) sont des immersions fermées et \(f\) un morphisme lisse, et l'on doit montrer que « \(i_*(E)\) \(n\)-pseudo-cohérent » équivaut à « \(j_*(E)\) \(n\)-pseudo-cohérent ». La question étant locale dans \(Y\) et dans \(Z\) au voisinage d'un point de \(X\), on peut supposer que \(X,Y,Z\) sont affines et qu'on a une factorisation
\[
\begin{tikzcd}
Y \arrow[r,hook,"i'"] \arrow[d,"f"'] & Z[T_1,\ldots,T_r]\arrow[dl,"f'"]\\
Z{,}&
\end{tikzcd}
\]
où \(i'\) est une immersion fermée, et la flèche oblique la projection canonique. Alors (VII 1.10) \(i'\) est une immersion régulière, donc (1.1.1) un objet \(F\) de \(D(Y)\) est \(n\)-pseudo-cohérent si et seulement si \(i'_*(F)\) l'est.

